\documentclass[12pt,a4paper]{report}
\usepackage[spanish]{babel}
\usepackage[document]{ragged2e} 
\addto\captionsspanish{\renewcommand{\contentsname}{Índice}}

\usepackage{titlesec}

% Modifica el formato del comando \chapter
\titleformat{\chapter}[hang]
{\normalfont\huge\bfseries} % Estilo del texto (Grande y Negrita)
{\thechapter\ - }           % Lo que va antes del título: "1 - "
{0pt}                       % Espacio entre el número y el título
{\MakeUppercase}            % Fuerza el título a estar en MAYÚSCULAS

\usepackage[margin=2cm]{geometry}
\usepackage{booktabs,tabularx,longtable}
\usepackage[table]{xcolor}
\usepackage{enumitem,fancyhdr,hyperref}
\usepackage{graphicx}
\usepackage{ragged2e}

\definecolor{ink}{HTML}{263746}\definecolor{red}{HTML}{9A3E3E}\definecolor{pale}{HTML}{FAEEEE}
\hypersetup{colorlinks=true,linkcolor=red,urlcolor=red}
\renewcommand{\headrulewidth}{0.4pt}

\begin{document}
\RaggedRight
\begin{titlepage}\centering{\color{ink}\Large\bfseries UNIVERSIDAD TECNOLÓGICA NACIONAL\par}\vspace{0.5cm}{\normalsize\bfseries INSTITUTO NACIONAL SUPERIOR DEL PROFESORADO TÉCNICO\par}

\vspace{1.6cm}\begin{figure}[h]
	\centering
	\includegraphics[width=0.35\textwidth]{imagenes/logoUTN.png}
\end{figure}
	
\vspace{1.3cm}{\Large Informe Final\par}
\vspace{0.1cm}{\selectfont\huge\bfseries Clickfud\par}
\vspace{0.1cm}{\selectfont\huge\bfseries Sistema  de pedidos de comida a domicilio ó para llevar\par}
\vspace{1.2cm}{\Large Trabajo final de Seminario 3.603}
\vspace{1.2cm}\textit{\Large \\Imperi, Jonatan Eduardo \\151957 \\jonatan.imperi@alu.inspt.utn.edu.ar}
\vspace{1.2cm}{\Large \\Fecha de presentación \\ \textit{09/11/2026}}
\vspace{1.2cm}{\Large{\\Docente de cátedra\\Prof. Mermoz, Daniel}}
\end{titlepage}

\thispagestyle{empty} % Quita el número solo en esta página

\newpage
~
\pagenumbering{Roman} % Cambia a romanos mayúsculas
\newpage

\section*{\Huge RESUMEN}
\begin{justify}
El desarrollo de ClickFud  tiene como finalidad entregar al cliente un sistema donde recibir y gestionar pedidos de su restaurante. Los consumidores tienen acceso a la página web por medio de cualquier navegador en donde hacen los pedidos, evitando de esta manera tener la obligación de descargar una aplicación dedicada a tal fin. El sistema completo se compone de tres interfaces visuales con las que se interactúan los distintos tipo de usuario.

La interfaz principal, utilizada por el administrador del local es una aplicación web que recibe y gestiona los pedidos entrantes, tener el control de las comandas y los pagos, tener el contacto de los clientes y además poseer un chat con el cual poder comunicarse con los repartidores a domicilio.

La segunda interfaz es una aplicación móvil para Android destinada a los repartidores del local para dar aviso de los envíos pendientes, confirmar las entregas y brindarles servicio de mapa con la ruta hacia el destino de entrega como también un chat para poder comunicarse con el restaurante.

La tercera interfaz orientada al público, es una tienda virtual que muestra las platos disponibles, ofrece a opción de retirar por el local o recibir el pedido a domicilio. Elegir el medio de pago y concretar la compra.
\end{justify}
\newpage
~
\newpage

\section*{\Huge ABSTRACT}
\begin{justify}
The development of ClickFud aims to provide the client with a system for receiving and managing restaurant orders. Consumers access the website via any web browser to place orders, thereby avoiding the need to download a dedicated app. The complete system consists of three interactive visual interfaces.

The main interface, used by the establishment's administrator, is a web application that receives and manages incoming orders; it handles order tracking and payments, maintains customedr contact details, and includes a chat feature for communicating with delivery drivers.

The second interface is an Android mobile app designed for the establishment's delivery drivers; it notifies them of pending deliveries, allows them to confirm completed deliveries, provides mapping services with routes to the destination, and includes a chat function for communicating with the restaurant.

The third interface, aimed at the public, is an online store that displays available dishes and offers the choice between pickup and home delivery, allowing customers to select a payment method and complete their purchase.
\end{justify}
\newpage
~
\newpage

\tableofcontents

\newpage
~
\newpage
\texttt{}
\chapter{Introducción}
\begin{justify}
El rubro gastronómico es de los que funcionan casi todos los días de la semana ya que la demanda de consumo lo requiere y hace posible. De esta manera, los locales que trabajan haciendo entregas a domicilio como también los que preparan el pedido para ser retirados por el cliente en el lugar, necesitan gestionar y procesar tales encargos. Hasta no hace mucho tiempo, los pedidos realizados de forma remota se hacían mediante llamadas telefónicas exclusivamente, y aunque hoy en día algunos optan por la mensajería instantánea, lo más habitual es el uso de aplicaciones que conectan el negocio de comida con miles de clientes cercanos y se encarga de la logística de entrega.
El principal problema que presentan dichas aplicaciones es el encarecimiento del producto final para el cliente como consecuencia del costo de las comisiones para los vendedores, lo cual puede inhibir el consumo.
El motivo para la creación de este sistema es reducir esos costos como también prescindir de sus métodos de venta que consisten en incentivar constantemente al consumidor con ofertas que duran muy poco tiempo.
\end{justify}

\pagenumbering{arabic}
	
\chapter{Estado del arte}
\begin{justify}
Las plataformas de delivery (pedidosYa, Rappi) ofrecen alcance y logística, pero cobran comisiones por pedido que reducen el margen del comercio. Por otro lado, los sistemas de gestión (FUDO, Bistrosoft, Kyte) resuelven la operación interna del local con un alcance funcional amplio, aunque no están orientados a un circuito de delivery propio de punta a punta. Este proyecto propone un modelo híbrido: un sistema de alcance acotado que integra en una sola solución la gestión de pedidos y cobros del local, una web de compra para clientes y una aplicación para repartidores propios, eliminando la comisión por pedido a costa de un menor número de funcionalidades.
\end{justify}

\section{PedidosYa}
\begin{justify}
Se encuentra entre las aplicaciones más conocidas de entrega a domicilio. PedidosYa no cobra una comisión fija única para todos, sino un porcentaje variable sobre el valor de cada pedido. En muchos casos arranca con promociones iniciales (por ejemplo, 5 por ciento el primer mes, 15 por ciento el segundo y un 25 por ciento a 30 por ciento a partir del tercer mes o según el tipo de contrato). Estos costos de comisión impulsan a que los locales aumenten sus precios finales en la aplicación en comparación a los precios de manera directa.
El esquema de pago hacia el negocio se hace de forma semanal
\end{justify}

\section{Rappi}
\begin{justify}
Muy parecida a PedidosYa pero con diferencias en condiciones de negocio, otro tipos de comisiones, distintos acuerdos, etc.
\end{justify}

\section{FUDO}
\begin{justify}
Fudo es una plataforma de gestión gastronómica en la nube que funciona 100 por ciento desde el navegador de una computadora, tablet o celular, sin necesidad de instalaciones complejas. Centraliza el salón, la cocina, el delivery, el control de stock, la caja y la facturación electrónica en un solo sistema.
\end{justify}
\chapter{Definición del problema}
La problemática ligada al estado del arte en donde las aplicaciones existentes tienen diferente formas de negocio, estas no cubren las necesidades mas elementales en su totalidad por si solas, mientras una tiene la logistica de entrega bien solucionada, traslada los precios finales de os productos. Y los que trabajan como gestión de negocio no resuelven de forma autónoma los envíos, ya que terminan tercer izando con las antes mencionadas.
Como dos pendientes a resolver son: Una solución propia de los envíos a domicilio y un sistema sencillo de manejar que cuente con lo necesario para el trabajo del día a día.
\chapter{Solución propuesta}

\section{Análisis de requisitos}

\section{Especificaciones del sistema}

\section{Análisis del sistema}

\chapter{Resultados}

\chapter{Conclusiones}

\section{Referencias}

\section{Bibliografía}

\section{Anexos}







\end{document}
